\begin{frame}{Überblick Vortrag}
\small
\begin{itemize}\setlength{\itemsep}{2mm}
  \item \textbf{\underline{Einführung und Ziele:}}
  \item \textbf{1. Rahmeninfos:} Pinecil, Temperatur, Spitzenpflege, Spitzenaktivator
  \item \textbf{2. Sicherheit:} Hitze, Arbeitsplatz, Dämpfe, Schutz, Verhalten, Erste Hilfe
  \item \textbf{3. Kurze Geschichte:} Antike, Elektronik
  \item \textbf{4. Wozu und Arten:} Was ist Löten, Anwendungen, Lötarten, THT, SMD
  \item \textbf{5. Werkzeuge und Material:} Lötstation, Spitzen, Lötzinn, Flussmittel, Entlöten, Helfer
\end{itemize}
\end{frame}

\begin{frame}{Überblick Vortrag}
\begin{itemize}\setlength{\itemsep}{2mm}
  \item \textbf{6. Techniken und Methoden:} Gute Lötstelle, Vorbereitung, Wärmeführung, Zinn zuführen, Absetzen, Nacharbeit, Temperatur und Zeit
  \item \textbf{7. Häufige Fehler und Lösungen:} Kalte Lötstelle, zu viel, zu wenig, Brücke, Überhitzung, Bauteil schief
  \item \textbf{8. Kontrolle und Praxis:} Inspektion, elektrische Prüfung, Ablauf Praxis, Merkhilfe
  \item \textbf{9. Projekte nach Erfahrungsstufe:} Stufen 1 bis 5, Sicherheit, Material, Vorstellung
  \item \textbf{Q \& A und Abschluss}
\end{itemize}
\end{frame}
